\documentclass[10pt,a4paper]{amsart}
\usepackage[margin=19mm]{geometry}
\usepackage{amsmath,amssymb,mathtools}
\usepackage[scr=rsfs,cal=euler]{mathalfa}
\usepackage{xfrac}
\usepackage[unicode,hidelinks,bookmarks=false]{hyperref}
\newcommand{\ie}{i.e.\ }
\newcommand{\cf}{cf.\ }
\newcommand{\resp}{respectively\ }
\newcommand{\Subsection}{\subsection}
\providecommand{\nobreakdash}{\nobreak\hbox{-}}
\setlength{\emergencystretch}{2em}
\multlinegap0pt
\usepackage{amssymb}
\usepackage{mathtools}
\usepackage[shortlabels]{enumitem}
\usepackage{bm}
\newtheorem{definition}{Definition}
\newtheorem{lemma}{Lemma}
\newtheorem{theorem}{Theorem}
\newcommand{\R}{\mathbb{R}}
\newcommand{\argmin}{\operatorname*{arg\,min}}
\newcommand{\hide}[1]{}
\newcommand {\la}   {{\langle}}
\newcommand {\ra}   {{\rangle}}
\title{Essential Simplices Dominate in Harmonic Representatives of One-Dimensional Persistent Classes}
\author{Saugata Basu, Aldo Guzm\'an-S\'aenz, Laxmi Parida}
\date{Source: arXiv:2607.26378v1 (CC BY 4.0)}
\begin{document}
\maketitle

\noindent\emph{Source and licence.} Essential Simplices Dominate in Harmonic Representatives of One-Dimensional Persistent Classes. Version arXiv:2607.26378v1 (CC BY 4.0), distributed by arXiv under the Creative Commons Attribution 4.0 International licence, http://creativecommons.org/licenses/by/4.0/, as recorded in the arXiv licence metadata for this exact version and shown on the abstract page https://arxiv.org/abs/2607.26378v1. Full licence notice: \texttt{licenses/arxiv-2607.26378-CC-BY-4.0.txt}.\par\smallskip

\noindent\textit{Editorial scope.} Counted unit: the article's theorem thm:simplified (source0.tex lines 1838--1860). The article's complete original proof of it is delivered with it (source0.tex lines 2184--2524, one \texttt{proof} environment of 341 lines, which the article places away from the statement). No mathematics is written, repaired or re-derived: the statement and the proof are the article's own lines, in the article's own notation, and no retained line is substituted or re-flowed.  Retained uncounted as the source-local context the proof needs: lines 1463--1481; lines 1484--1496; lines 1538--1543; lines 1899--1926; lines 2138--2145.  Nothing was omitted from any retained span, so the package carries no omission record.  No retained line contains a citation, so the wrapper carries no bibliography and no bibliography entry.  No retained line contains a figure environment, an image-inclusion command or drawing code, so no image is delivered and the package declares no asset dependency.  Editorial formatting only, in addition: an AMS \texttt{amsart} preamble that restates the article's own declarations and macros used by the retained lines, namely the environments \texttt{definition}, \texttt{lemma}, \texttt{theorem} on separate counters, together with the standard packages those lines need.  The retained environments print in this document as Definition 1, Lemma 1, Theorem 1 in order of appearance, whereas the article prints the same items with the numbering of its own full document.  The complete native source of the article as distributed in the arXiv e-print for this exact version is bundled unchanged as \texttt{sources/206265/source0.tex}.
\begin{definition}[Relative essential simplex]
\label{def:relative-essential}
We call a simplex \(\sigma\) of \(K^{[p]}\) \emph{relatively essential}
with respect to \(h\) if
\[
\langle h+w,\sigma\rangle\neq 0
\qquad
\text{for every }w\in W_p.
\]
Equivalently,
\[
\langle h+b+z',\sigma\rangle\neq 0
\qquad
\text{for every }b\in B_p(K),\ z'\in Z_p(K').
\]

We denote by  \(S_E(h)\)  the set of relatively essential simplices with respect to $h$, and by
\(S_N(h) = K^{[p]} \setminus S_E(h)\). 
\end{definition}

\begin{lemma}[Alternative characterization of relative essentiality]
\label{lem:linear}
For \(\sigma \in K^{[p]}\), the following are equivalent (see Definition~\ref{def:relative-essential}):
\[
\sigma\in S_E(h)
\]
and
\[
\langle h,\sigma\rangle\neq 0
\quad\text{and}\quad
\sigma\perp W_p.
\]
\end{lemma}

\begin{lemma}
\label{lem:orthogonality}
\[
h\perp W_p.
\]
\end{lemma}

\begin{lemma}[Minimum-norm prescribed-boundary chains]
\label{lem:graph}
Let \(G\) be a finite oriented graph. Let \(\sigma\) be an oriented edge of
\(G\), with
$
\partial\sigma=b-a
$,
and let
$
G'=G\setminus\{\sigma\}
$.
Suppose 
\(x\in C_1(G')\) is the unique minimum-norm solution of
$
\partial_{G'} x=a-b
$,
%%sb adds
assuming the solution set is non-empty.
Then, for every edge \(\tau\) of \(G'\) that is not a bridge separating \(a\)
from \(b\),
$
|\langle x,\tau\rangle|<1
$.
In particular, if \(\tau\) lies in a cycle of \(G'\), then
$
|\langle x,\tau\rangle|<1
$.
\end{lemma}

\begin{lemma}[Cycle coefficient implies non-bridge]
\label{lem:bridge}
Let \(G\) be a finite graph, and let \(c\in Z_1(G)\). If
\[
\langle c,\tau\rangle\neq 0
\]
for an edge \(\tau\), then \(\tau\) is not a bridge of \(G\).
\end{lemma}

\begin{theorem}
\label{thm:simplified}
Let \(K\) be a finite simplicial complex and \(K'\subseteq K\) a subcomplex.
Assume
\begin{equation}
\label{eqn:thm:simplified}
\dim_{\mathbb R} H_1(K)/i_1(H_1(K'))=1.
\end{equation}
Let $h$ be a non-zero element of \(\mathcal H_1(K)\) spanning
\(\mathcal H_1(K)\cap \mathfrak{i}_1(\mathcal H_1(K'))^\perp\).
%%sb rewrites
\hide{
Then for every \(\sigma\in S_E(h)\) and every \(\tau\in S_N(h)\),
\[
|\langle h,\tau\rangle|<|\langle h,\sigma\rangle|.
\]
}
Then for every \(\sigma\in S_E(h)\) and every \(\tau\in K^{[1]}, \tau \neq \sigma\),
\[
|\langle h,\tau\rangle| \leq |\langle h,\sigma\rangle|,
\]
with equality if and only if $\tau \in S_E(h)$.
\end{theorem}

\begin{proof}[Proof of Theorem~\ref{thm:simplified}]
Fix
\[
\sigma\in S_E(h).
\]
By the characterization of relative essentiality (Lemma~\ref{lem:linear}),
\[
\langle h,\sigma\rangle\neq 0
\qquad\text{and}\qquad
\sigma\perp W_1.
\]

Set
\[
h_\sigma:=\langle h,\sigma\rangle,
y:=\frac{h}{h_\sigma}.
\]
Then
$
\langle y,\sigma\rangle=1
$.

Let
$
\ell_\sigma:Z_1(K)\to\R
$
be the linear functional defined by 
\[
\ell_\sigma(z)=\langle z,\sigma\rangle.
\]

Since
$
\sigma\perp W_1
$,
the functional \(\ell_\sigma\) vanishes on \(W_1\). Hence it descends to a
linear functional
$
\overline{\ell}_\sigma:
Z_1(K)/W_1
\to
\R
$.


Moreover,
\[
\overline{\ell}_\sigma([h])
=
\langle h,\sigma\rangle
=
h_\sigma
\neq 0.
\]

It follows from the hypothesis \eqref{eqn:thm:simplified} that
$
\dim Z_1(K)/W_1=1
$,
and hence 
the functional \(\overline{\ell}_\sigma\) is injective. Therefore
$
\ker\ell_\sigma=W_1.
$

Equivalently,
\[
W_1
=
\{z\in Z_1(K):\langle z,\sigma\rangle=0\}.
\]

Define the affine space
\[
\mathcal A_\sigma
=
\{z\in Z_1(K):\langle z,\sigma\rangle=1\}.
\]

Since
$
\langle y,\sigma\rangle=1
$,
we have
$
y\in\mathcal A_\sigma.
$

If
$
z\in\mathcal A_\sigma
$,
then
$
\langle z-y,\sigma\rangle=0
$
and
$
z-y\in Z_1(K)
$.
Therefore
$
z-y\in W_1
$.
Hence
$
\mathcal A_\sigma=y+W_1
$.

%%By the previous orthogonality lemma,
Using Lemma~\ref{lem:orthogonality},
$
h\perp W_1
$,
which implies that 
$
y\perp W_1
$.
Therefore \(y\) is the unique minimum-norm element of \(\mathcal A_\sigma\):
\[
y
=
\argmin
\left\{
\|z\|^2:
z\in Z_1(K),\ \langle z,\sigma\rangle=1
\right\}.
\]

Let
$
G=K^{(1)}
$,
%%be the \(1\)-skeleton of \(K\), 
and let
$
G'=G\setminus\{\sigma\}
$.

Let
$
\partial\sigma=b-a
$.

For any
$
z\in\mathcal A_\sigma
$,
the chain
$
z-\sigma
$
is supported on \(G'\), and
\[
\partial(z-\sigma)
=
-\partial\sigma
=
a-b.
\]

Since the coefficient of \(\sigma\) is fixed to be \(1\), minimizing
$
\|z\|^2
$
over \(z\in\mathcal A_\sigma\) is equivalent to minimizing
$
\|z-\sigma\|^2
$
among chains on \(G'\) with boundary \(a-b\).

Thus
$
x:=y-\sigma
$
is the unique minimum-norm chain in \(C_1(G')\) satisfying
$
\partial x=a-b
$.

We now break the proof into two cases.
\begin{enumerate}
\item Case $\tau\in S_N(h)$.

There are now two sub-cases.
\begin{enumerate}
\item
Suppose
\[
\langle h,\tau\rangle=0.
\]
Then
\[
|\langle h,\tau\rangle|
=
0
<
|\langle h,\sigma\rangle|,
\]
because \(\sigma\in S_E(h)\) implies
$
\langle h,\sigma\rangle\neq 0
$.

\item
Now suppose 
$
\langle h,\tau\rangle\neq 0.
$

Since
$
\tau\in S_N(h)
$,
the edge \(\tau\) is not relatively essential. Because its \(h\)-coefficient is
nonzero, 
%%the characterization of relative essentiality 
Lemma~\ref{lem:linear}
implies that
$
\tau\not\perp W_1
$.
Thus there exists
$
w\in W_1
$
such that
$
\langle w,\tau\rangle\neq 0
$.

Since
$
\sigma\in S_E(h),
$
we have
$
\sigma\perp W_1.
$
Therefore
$
\langle w,\sigma\rangle=0.
$

Also,
$
w\in W_1\subseteq Z_1(K)
$,
and 
so \(w\) is a cycle in \(K\). Since its \(\sigma\)-coefficient is zero, \(w\)
is a cycle in \(G'\). Since
$
\langle w,\tau\rangle\neq 0
$,
%%the previous bridge lemma 
Lemma~\ref{lem:bridge}
implies that \(\tau\) is not a bridge of \(G'\).

By 
%%the minimum-norm prescribed-boundary lemma 
Lemma~\ref{lem:graph}
applied to
$
x=y-\sigma
$,
we get
$
|\langle x,\tau\rangle|<1
$.

Since \(\tau\neq\sigma\), and \(x=y-\sigma\), we have
\[
\langle x,\tau\rangle
=
\langle y,\tau\rangle.
\]
Hence
$
|\langle y,\tau\rangle|<1
$.

Finally,
$
y=\frac{h}{\langle h,\sigma\rangle}
$, and
hence
\[
\langle y,\tau\rangle
=
\frac{\langle h,\tau\rangle}{\langle h,\sigma\rangle}.
\]
Thus
\[
\left|
\frac{\langle h,\tau\rangle}
     {\langle h,\sigma\rangle}
\right|
<1.
\]
Equivalently,
\[
|\langle h,\tau\rangle|
<
|\langle h,\sigma\rangle|.
\]
\end{enumerate}
The two cases exhaust all
$
\tau\in S_N(h)
$.

\item Case $\tau \in S_E(h)$. 
Suppose that $z \in Z_1(K)$. We claim that 
\begin{equation}
\label{eqn:proof:thm:simplified:0}
  \la z, \sigma\ra \neq 0  \Rightarrow \la z, \tau\ra \neq 0
\end{equation}
Indeed $\tau \perp W_1$ and $\la h, \tau \ra \neq 0$. Moreover,  
$Z_1(K) = \operatorname{span}(h) \oplus W_1$, and hence
$z = c h + w$ for some $c \in \mathbb{R}, w \in W_1$. Since $\sigma \perp W_1$,
$c \neq 0$. This implies that 
$\la z , \tau\ra = c \cdot \la h,\tau\ra \neq 0$ as well. 

Since $\operatorname{ker}(\ell_\sigma) = \operatorname{ker}(\ell_\tau) = W_1$,
and $W_1$ is of codimension $1$ in $Z_1(K)$, the linear functional 
\begin{equation}
\label{eqn:proof:thm:simplified}
\ell_\sigma = \lambda \ell_\tau,
\end{equation}
for some $\lambda \neq 0$. 
Since $\la h, \sigma \ra \neq 0$ and $h \in Z_1(K)$, Lemma~\ref{lem:bridge} implies that $\sigma$ is not a bridge in $G$. So there exists a simple cycle $z_1$
(i.e.
$z_1 \in Z_1(K)$ with all coefficients in $\{0,1,-1\}$), with $\sigma$ in its support. Then 
$\ell_\sigma(z_1) = \la z_1,\sigma\ra = \pm 1$, and hence $\ell_\tau(z_1) = \la z_1,\tau\ra  \neq 0$ using \eqref{eqn:proof:thm:simplified:0}. Therefore,
$\ell_\tau(z_1) \in \{1,-1\}$. 
Comparing with \eqref{eqn:proof:thm:simplified} we get that $\lambda \in \{1,-1\}$.
Evaluating $\ell_\sigma$ and $\ell_\tau$ at $h$, we obtain
$\la h, \tau\ra = \pm \la h, \sigma\ra$.
\end{enumerate}
Therefore the theorem is proved.
\end{proof}

\end{document}
